\section{Related Work}
\label{sec:related}

\paragraph{LLM judges as reward.} Replacing human preference labels with model-generated ones
\citep{christiano2017deeprl,ouyang2022instructgpt} is now routine, via AI feedback
\citep{guo2024oaif}, self-rewarding loops \citep{yuan2024selfrewarding}, or best-of-$N$
distillation \citep{pace2024westofn}. Such judges carry systematic biases --- position and
verbosity \citep{zheng2023judging}, length \citep{dubois2024lengthcontrolled}, self-preference
\citep{panickssery2024selfpreference}. Our setting inherits all of these and adds a coupling:
the same model family simulates the patient \emph{and} grades the transcript, so the policy
optimises a signal produced by its own conversational partner. The held-out judge in
\S\ref{sec:heldout} controls for exactly that.

\paragraph{Reward hacking and overoptimisation.} Optimising a proxy past the point where it
tracks the objective is well characterised theoretically \citep{skalse2022defining} and
empirically for learned reward models \citep{gao2023scaling}. Sycophancy --- agreement and
praise that raise ratings without improving the response --- is the failure mode closest to
what we observe \citep{sharma2024sycophancy,perez2023discovering}. What is usually missing is
a cheap way to \emph{detect} it during a training run without new human labels. Gain
retention (\S\ref{sec:heldout}) is our attempt at one; it needs a second grader and no
annotation, and it separates the two arms four iterations before their outcome curves
diverge.

\paragraph{Tree-structured preference data.} Constructing preferences by branching a dialogue
and scoring simulated continuations connects to MCTS-guided preference learning
\citep{xie2024mctsdpo,yu2023promptmcts,yuan2024eurus} and to the preference-tree formulation
this work builds on \citep{baruch2025pto}, which also introduced a look-ahead depth $K$
(candidates scored after $K$ further simulated turns). We hold $K=0$ throughout so the
optimizer is the only variable; the $K$ comparison is reported separately.

\paragraph{MI and computational models of counselling.} MI's coding systems
\citep{moyers2016miti} give this domain what most dialogue benchmarks lack: an externally
specified account of which behaviours are correct, with published competency thresholds.
LLM-simulated patients have been used to train and assess counsellors
\citep{steenstra2024scaffolding}, \citet{yosef2024assessing} validate the questionnaire-filling
evaluators we use as the reward, and \citet{chen2025broaden} study LLM therapist behaviour
beyond aggregate quality. Ours is not a new MI system but a measurement result about what
happens when such an evaluator is optimised against for long enough.
